\documentclass[11pt]{article}

\usepackage[review]{acl}

\usepackage{times}
\usepackage{latexsym}
\usepackage[T1]{fontenc}
\usepackage[utf8]{inputenc}
\usepackage{microtype}
\usepackage{inconsolata}
\usepackage{graphicx}
\usepackage{booktabs}
\usepackage{caption}
\usepackage{amsmath}

\title{Probing Functional Correctness in Diffusion Language Models}

\author{First Author \\
  Affiliation \\
  \texttt{email@domain} \\\And
  Second Author \\
  Affiliation \\
  \texttt{email@domain} \\}

\begin{document}
\maketitle

\begin{abstract}
Diffusion language models generate text by iteratively denoising all tokens in parallel, but when and where their hidden states encode whether the output will be functionally correct remains unknown.
We present the first probing study of diffusion LLM internals, training linear classifiers on hidden states extracted at intermediate denoising steps to predict functional correctness.
Across two models (LLaDA-8B, Dream-7B) and two tasks (JSON schema validation, math reasoning), we find that (1) correctness signal is already present at step~0, before any denoising occurs, (2) signal emergence is task-dependent, with structural tasks showing flat profiles and reasoning tasks showing gradual buildup, and (3) layer dynamics diverge across architectures.
We further demonstrate that per-step probe confidence can identify instances likely to fail, enabling selective generation that avoids 36--98\% of wasted compute depending on the task.
\end{abstract}

\section{Introduction}

Diffusion language models (dLLMs) have emerged as a promising alternative to autoregressive (AR) generation \citep{austin2021d3pm,sahoo2024simple,lou2024discrete}.
Unlike AR models, which generate tokens sequentially, dLLMs produce all output tokens simultaneously through iterative denoising: starting from a fully masked or noisy sequence and progressively refining it over many steps.

This parallel generation process raises a fundamental question: \emph{when} during denoising, and \emph{where} within the model's layers, does the network encode whether its output will be functionally correct?
In AR models, probing classifiers have revealed how linguistic and factual knowledge localizes across layers and positions \citep{tenney2019bert,geva2023dissecting,azaria2023internal}.
However, no analogous study exists for diffusion LLMs, whose bidirectional, iterative generation process creates a fundamentally different representational landscape.

We address this gap with the first probing study of dLLM hidden states.
Our contributions are: (1) we show that correctness signal exists from step~0, exposed simultaneously across all generation positions, a structural advantage of dLLMs over AR models where such signal is concentrated in the final prompt token; (2) we identify task-dependent emergence patterns, with structural tasks decided early and reasoning tasks requiring iterative refinement; and (3) we reveal divergent layer dynamics across architectures, with LLaDA concentrating signal in upper layers while Dream migrates signal from upper to lower layers during denoising.

\section{Background}

\paragraph{Diffusion language models.}
LLaDA \citep{nie2025llada} uses block masking during training and generates by progressively unmasking tokens over a fixed schedule.
Dream \citep{ye2025dream} employs a global masking schedule with a different unmasking strategy.
Both models condition on a fully visible prompt and iteratively denoise a generation region of masked tokens.

\paragraph{Probing classifiers.}
Probing involves training simple classifiers on frozen model representations to test whether specific information is encoded \citep{belinkov2022probing,conneau2018cram}.
We adopt this methodology for dLLMs, probing hidden states at intermediate denoising steps.

\paragraph{Related work.}
Prior studies have probed AR model representations for linguistic structure \citep{tenney2019bert}, factual knowledge \citep{geva2023dissecting}, and truthfulness \citep{azaria2023internal}.
Confidence-based methods for AR models use early-exit thresholds to skip computation \citep{schuster2022confident,elbayad2020depth}.
Two concurrent works study dLLM representations from complementary angles.
\citet{shnaidman2025steering} demonstrate activation steering for safety behaviors in masked diffusion LLMs, showing that dLLM representations are amenable to linear intervention.
\citet{wang2025time} identify temporal oscillation in dLLM outputs, where correct answers appear mid-denoising but are overwritten in later steps.
\citet{mundler2025constrained} propose constrained decoding with context-free grammars to enforce syntactic correctness in dLLM outputs; our work instead asks whether the model's own representations already encode whether the output will be correct, without external constraints.
We complement these prior findings by probing hidden representations for functional correctness across the full layer-step grid, characterizing \emph{when} and \emph{where} correctness signal emerges.

\section{Method}

\paragraph{Hidden state extraction.}
Given a dLLM with $L$ layers, we run 128-step denoising and extract hidden states at 7 checkpoint steps $t \in \{0, 1, 4, 16, 32, 64, 127\}$.
At each checkpoint, we obtain the full hidden state tensor of shape $(L, \text{gen\_length}, d)$.
We partition the generation region into 4 equal-length position regions and mean-pool within each region, yielding features of shape $(L, 4, d)$ per instance per step.
We choose 4 regions to align with LLaDA's block length of 32 tokens (gen\_length / block\_length = 4 for JSON schema's 128-token blocks) while keeping the number of features manageable given the limited sample sizes.
An ablation over 1, 2, and 4 regions shows best AUC varies by less than 0.02 across configurations (Appendix~\ref{sec:appendix_region}).

\paragraph{Probe training.}
For each (step, layer, region) triple, we reduce dimensionality with PCA to 64 components, apply standard scaling, and train a logistic regression classifier to predict functional correctness \citep{pedregosa2011scikit}.
We evaluate with 5-fold stratified cross-validation, reporting AUC as our primary metric.
Control probes trained on shuffled labels yield AUC $\approx 0.50$ across all settings, confirming that the observed signal is genuine and not an artifact of the probing pipeline.

\paragraph{Functional correctness.}
For JSON schema validation, an output is correct if it parses as valid JSON and matches the reference structure.
For math reasoning (GSM8K), an output is correct if the extracted numeric answer matches the ground truth.

\paragraph{Selective generation.}
We simulate instance-level filtering by training a per-step probe (using the best layer and region from the full analysis).
At inference time, if the probe predicts with confidence exceeding a threshold $\tau$ that an instance will fail, we skip the remaining denoising steps for that instance entirely, producing no output rather than wasting compute on a generation predicted to be incorrect.
We report the fraction of instances filtered and the probe's classification accuracy.

\section{Experimental Setup}

\paragraph{Models.}
We evaluate LLaDA-8B-Instruct (33 layers) \citep{nie2025llada} and Dream-7B-Instruct (29 layers) \citep{ye2025dream}, two open-weight diffusion LLMs with different masking strategies.

\paragraph{Datasets.}
We use two tasks that test different capabilities:
(1) \textbf{JSON schema validation} (272 instances from json-mode-eval-extended), a structural task requiring syntactically and semantically correct JSON output; and
(2) \textbf{GSM8K} (1,319 test instances) \citep{cobbe2021gsm8k}, a multi-step math reasoning task.
Generation lengths are 256 and 512 tokens, respectively.

\paragraph{Baseline functional rates.}
Table~\ref{tab:baselines} shows the functional correctness rates with 128 denoising steps.
Both datasets have reasonable class balance, enabling meaningful probe training.

\begin{table}[t]
  \centering
  \small
  \begin{tabular}{lcc}
    \toprule
    & \textbf{JSON Schema} & \textbf{GSM8K} \\
    \midrule
    LLaDA-8B & 48.5\% (132/272) & 66.3\% (875/1319) \\
    Dream-7B & 46.0\% (125/272) & 61.5\% (811/1319) \\
    \bottomrule
  \end{tabular}
  \caption{Baseline functional correctness rates (seed=0, 128 denoising steps).}
  \label{tab:baselines}
\end{table}

\section{Results}

\subsection{Task-Dependent Emergence}

Figure~\ref{fig:heatmap} shows AUC heatmaps across layers and denoising steps for all four model-task combinations.
Two distinct emergence patterns are visible.

For \textbf{JSON schema}, correctness signal is already strong at step~0 (AUC $\approx 0.78$--$0.80 \pm 0.04$) and remains relatively flat across all denoising steps.
This suggests that the prompt encoding alone, before any denoising occurs, already contains sufficient information to predict whether the output will be structurally correct.

For \textbf{GSM8K}, step-0 signal is already well above chance (AUC $\approx 0.67$--$0.71 \pm 0.03$), but denoising brings substantially more gain: AUC rises to $\approx 0.78$--$0.82 \pm 0.03$ by step~127 (Figure~\ref{fig:auc_curve}).

Both tasks have substantial step-0 signal, but JSON schema gains only $\Delta\text{AUC} \approx 0.04$ from denoising while GSM8K gains $\approx 0.11$, reflecting the difference between structural and reasoning tasks.
Standard deviations across 5-fold CV are $\leq 0.05$ for all configurations, indicating stable estimates.
Unlike AR models where such signal is concentrated at the final prompt token, dLLMs expose correctness-relevant information simultaneously across all generation positions via bidirectional attention.

Seed reranking at step~1 yields +0.0\% improvement for both models, confirming that the probe captures instance-level difficulty rather than seed-specific quality.

\begin{figure}[t]
  \centering
  \includegraphics[width=\columnwidth]{fig1_heatmap.png}
  \caption{AUC heatmaps (layer $\times$ step) for predicting functional correctness. Stars mark the best layer per step. JSON schema shows strong signal from step~0, while GSM8K shows gradual emergence. Note the layer migration pattern in Dream (top-right, bottom-right).}
  \label{fig:heatmap}
\end{figure}

\begin{figure}[t]
  \centering
  \includegraphics[width=\columnwidth]{fig2_auc_curve.png}
  \caption{Best AUC across layers at each denoising step. JSON schema is flat ($\approx$0.80 from step~0), while GSM8K rises from $\approx$0.71 to $\approx$0.82.}
  \label{fig:auc_curve}
\end{figure}

\subsection{Divergent Layer Dynamics}

The heatmaps in Figure~\ref{fig:heatmap} reveal strikingly different layer dynamics between the two architectures.
Table~\ref{tab:best_layer} shows the best layer (by AUC) at each denoising step.

\textbf{LLaDA} concentrates its best probing signal consistently in upper layers (L22--28) across all denoising steps and both tasks.
This is similar to patterns observed in AR transformers, where upper layers encode more task-relevant features.

\textbf{Dream} shows a task-dependent migration pattern.
On JSON schema, the best layer migrates from the upper range (L17--23 at steps 0--16) to lower layers (L4--7 at steps 32--127).
On GSM8K, however, Dream's best layers remain in the upper range (L19--23) throughout denoising, with no downward migration.
This suggests that the migration occurs only when the task is simple enough for the model to resolve early.

\begin{table}[t]
  \centering
  \small
  \begin{tabular}{l*{7}{c}}
    \toprule
    \textbf{Step} & 0 & 1 & 4 & 16 & 32 & 64 & 127 \\
    \midrule
    \multicolumn{8}{l}{\textit{JSON Schema}} \\
    ~~LLaDA & 22 & 25 & 26 & 25 & 15 & 25 & 24 \\
    ~~Dream & 17 & 21 & 23 & 19 & \textbf{4} & \textbf{7} & \textbf{5} \\
    \midrule
    \multicolumn{8}{l}{\textit{GSM8K}} \\
    ~~LLaDA & 24 & 24 & 24 & 25 & 28 & 24 & 25 \\
    ~~Dream & 23 & 21 & 21 & 21 & 21 & 21 & 19 \\
    \bottomrule
  \end{tabular}
  \caption{Best layer (by AUC) at each denoising step. Bold values highlight Dream's migration to lower layers on JSON schema. No migration occurs on GSM8K.}
  \label{tab:best_layer}
\end{table}

We hypothesize that this difference relates to training origin: LLaDA is trained from scratch, while Dream is fine-tuned from the AR model Qwen-2.5 and may inherit its pretrained layer hierarchy.
The AR inductive bias may cause Dream to initially rely on upper layers but redistribute information to lower layers as diffusion-specific dynamics take over, a redistribution that only occurs when the task is simple enough.

\subsection{Selective Generation}

We evaluate whether per-step probe confidence can identify instances whose outputs will be incorrect, enabling the system to skip generation entirely for hopeless instances.
Table~\ref{tab:early_exit} shows results with a confidence threshold of $\tau = 0.80$.

\begin{table}[t]
  \centering
  \small
  \begin{tabular}{lcc}
    \toprule
    & \textbf{JSON Schema} & \textbf{GSM8K} \\
    \midrule
    \multicolumn{3}{l}{\textit{LLaDA-8B}} \\
    ~~Compute saved & 98.4\% & 61.3\% \\
    ~~Clf.\ accuracy & 73.5\% & 73.5\% \\
    \midrule
    \multicolumn{3}{l}{\textit{Dream-7B}} \\
    ~~Compute saved & 96.7\% & 36.0\% \\
    ~~Clf.\ accuracy & 72.4\% & 73.5\% \\
    \bottomrule
  \end{tabular}
  \caption{Selective generation results ($\tau = 0.80$). Compute saved is the fraction of denoising steps avoided by skipping instances predicted to fail (no output is produced for skipped instances). Clf.\ acc is the fraction of instances for which the probe correctly predicts functional vs.\ non-functional.}
  \label{tab:early_exit}
\end{table}

The results are strongly task-dependent.
For JSON schema, both models save over 96\% of compute: because the probe is confident from step~0, it identifies most instances predicted to fail immediately and skips their generation entirely.
For GSM8K, savings are more modest (36--61\%) because reasoning tasks require more denoising steps before the probe can confidently distinguish hopeless instances.
In all cases, the probe's classification accuracy remains above 72\%, suggesting that it reliably identifies instances that will fail.

\section{Discussion}

\paragraph{Why does step-0 signal exist?}
At step~0, the dLLM has performed a forward pass over the prompt with masked outputs, producing representations that reflect the model's ``assessment'' of the generation task.
For structural tasks like JSON generation, this assessment is highly predictive: the model may already encode whether the prompt is one it can handle correctly.
For reasoning tasks, the prompt assessment provides weaker signal, and the model must iteratively refine its representations through denoising.
The dLLM step-0 signal spans all generation positions simultaneously, arising from bidirectional attention over prompt and masked tokens.
An AR baseline probe on Qwen-2.5-7B's last prompt token achieves lower AUC (0.69--0.74) than dLLM step-0 probes (0.71--0.80) on the same instances (Appendix~\ref{sec:appendix_ar}), confirming that bidirectional attention over masked generation tokens provides signal beyond what AR prompt encoding captures.

\paragraph{Relationship to temporal oscillation.}
\citet{wang2025time} observe that dLLM outputs can oscillate during denoising: correct answers appear at intermediate steps but are overwritten later.
Our probe measures whether hidden states \emph{encode} correctness, not whether the current surface output is correct.
On JSON schema, hidden-state signal is flat from step~0 while surface tokens are still being refined, suggesting the model's internal plan is established early and oscillation reflects surface-level instability rather than changes in the underlying representation of correctness.

\paragraph{Layer dynamics and training origin.}
LLaDA (trained from scratch) shows stable upper-layer localization resembling AR models, while Dream (fine-tuned from Qwen-2.5) exhibits layer migration on simple tasks but not complex ones.
This suggests AR pretraining establishes an upper-layer bias that persists for difficult tasks, but diffusion fine-tuning enables redistribution to lower layers when the task is simple enough.

\paragraph{Practical implications.}
The selective generation results suggest that dLLMs waste significant compute on instances whose outputs will be incorrect.
A lightweight probe could enable instance-level filtering, skipping generation for hopeless instances and reserving compute for those likely to succeed.
The task dependence of our results suggests that structured generation tasks benefit most from aggressive filtering, while reasoning tasks require more conservative thresholds.

\section{Conclusion}

We present the first probing study of diffusion language model internals, revealing that functional correctness signal is encoded from the very first denoising step, with emergence patterns that depend on task complexity and layer dynamics that differ across architectures.
Concretely, step-0 probes achieve AUC 0.78--0.80 on JSON schema and 0.71 on GSM8K, signal peaks at intermediate steps (step~64) rather than the final step, and selective generation based on probe confidence can avoid 36--98\% of wasted compute depending on the task and model.
A comparison with an AR baseline (Qwen-2.5-7B) confirms that the dLLM step-0 signal is not merely prompt difficulty detectable by any model, but reflects additional information from bidirectional attention over masked generation tokens.
Length-stratified controls confirm that correctness signal persists (AUC 0.75--0.81) even after removing the length confound.

These findings open several directions: integrating probe-based selective generation into actual inference pipelines, extending the analysis to more dLLM architectures and tasks, and exploring nonlinear or multi-layer probes for adaptive computation.

\section*{Limitations}

Our probe is a simple logistic regression classifier with PCA-reduced features, which may underestimate the information available in hidden states.
Nonlinear probes or probes with access to multiple layers simultaneously might reveal richer structure.

The selective generation analysis is a simulation: we evaluate probe confidence on states extracted during full 128-step generation, but do not validate the approach in an actual inference pipeline where instances predicted to fail would be skipped entirely.
Integration would require running the probe at each denoising step and filtering instances in real time, introducing inference overhead that our simulation does not account for.

We study only two models (LLaDA-8B, Dream-7B) and two tasks (JSON schema validation, GSM8K).
Broader coverage across more architectures (e.g., MDLM, SEDD) and tasks (e.g., code generation, translation) is needed to establish generality.
Both models are in the 7--8B parameter range; it is unclear whether our findings hold for smaller or larger diffusion LLMs.

A control probe trained to predict reference output length (binary median split) achieves higher AUC (0.86--0.99) than the correctness probe (0.78--0.85) across all configurations (Appendix~\ref{sec:appendix_length}), indicating that hidden states encode output length more strongly than functional correctness.
To control for this confound, we re-train the correctness probe on a length-matched subsample where functional and non-functional instances have identical length distributions.
AUC drops modestly (0.75--0.81 vs.\ 0.79--0.85 on the full data) but remains well above chance, confirming that the probe captures genuine correctness signal beyond output length (Appendix~\ref{sec:appendix_length}).
Other surface-level confounders, such as prompt length or the presence of specific formatting tokens, may also inflate probe performance and are not controlled for in this work.

We extract hidden states at only 7 of the 128 denoising steps (log-spaced). This coarse temporal granularity may obscure fine-grained dynamics. For example, Dream-7B's best layer on JSON schema jumps from layer~19 at step~16 to layer~4 at step~32; without observing the intermediate steps, we cannot determine whether this is a gradual migration or a sudden phase transition.

Finally, probing is inherently correlational: linear separability of a feature does not guarantee that the model actively uses it during generation \citep{belinkov2022probing}.
Causal validation through activation interventions, verifying that perturbing the identified representations changes output correctness, remains an important direction for future work.

\bibliography{custom}

\appendix

\section{Generation and Probing Details}
\label{sec:appendix_generation}

\paragraph{Denoising procedure.}
Both models run 128 denoising steps.
LLaDA uses block-based unmasking: the generation region is divided into blocks of 32 tokens, and each block is denoised sequentially over $128 / (\text{gen\_length} / 32)$ steps per block.
Within each block, the model selects the most confident tokens to unmask at each step.
Dream uses a global linear schedule: at step $i$, the fraction of tokens to unmask is determined by a linearly decreasing timestep schedule from $t=1$ to $t=\epsilon$, and the most confident masked tokens are unmasked globally.

\paragraph{Prompt construction.}
For JSON schema, we prepend a system prompt containing the target schema and append \texttt{```json\textbackslash n} as a generation prefix.
For GSM8K, the system prompt instructs the model to solve step-by-step and end with \texttt{\#\#\#\#} followed by the numeric answer.
Both models use their respective chat templates.

\paragraph{Hidden state extraction.}
At each of the 7 checkpoint steps $\{0, 1, 4, 16, 32, 64, 127\}$, we run a forward pass with \texttt{output\_hidden\_states=True} and extract the full hidden state tensor across all layers.
The generation region is partitioned into 4 equal-length position regions, and we mean-pool within each region, producing a feature vector of dimension $d$ (4096 for LLaDA, 3584 for Dream) per layer per region per instance.

\paragraph{Probe pipeline.}
For each (step, layer, region) configuration, we construct the feature matrix $X \in \mathbb{R}^{n \times d}$ and apply:
(1) \texttt{StandardScaler} to zero-mean and unit-variance normalize features,
(2) PCA reduction to 64 components,
(3) logistic regression with L2 regularization ($C = 1.0$, LBFGS solver, max 1000 iterations).
We use 5-fold stratified cross-validation (random state 42) and report the mean AUC across folds.

\paragraph{Computational resources.}
All generation and feature extraction use NVIDIA A100 (80GB) GPUs on Modal, with 8 parallel A100s per experiment.
Models are loaded in bfloat16 precision.
Probe training runs on CPU.
Total compute: approximately 30 A100-hours across all experiments.

\section{Selective Generation Details}
\label{sec:appendix_early_exit}

The selective generation simulation works as follows.
For each instance, we iterate through the checkpoint steps in order.
At each step, we compute the probe's confidence as $\max(p, 1-p)$ where $p$ is the predicted probability of functional correctness.
If the probe confidently predicts failure (confidence exceeds threshold $\tau$ and $p < 0.5$), we skip the remaining denoising steps for that instance, producing no output.
If no checkpoint triggers skipping, the instance proceeds to full generation.
Compute savings are computed as $1 - \frac{t_{\text{skip}} + 1}{128}$, where $t_{\text{skip}}$ is the step at which filtering occurs.

The main paper reports results at $\tau = 0.80$.
Table~\ref{tab:early_exit_full} shows results across a range of thresholds for all configurations.

\begin{table}[t]
  \centering
  \small
  \begin{tabular}{lcccc}
    \toprule
    $\tau$ & \multicolumn{2}{c}{\textbf{JSON Schema}} & \multicolumn{2}{c}{\textbf{GSM8K}} \\
    & Saved & Clf.\ acc & Saved & Clf.\ acc \\
    \midrule
    \multicolumn{5}{l}{\textit{LLaDA-8B}} \\
    0.60 & 99.2\% & 64.0\% & 83.4\% & 67.2\% \\
    0.70 & 99.0\% & 69.5\% & 74.2\% & 70.3\% \\
    0.80 & 98.4\% & 73.5\% & 61.3\% & 73.5\% \\
    0.90 & 93.3\% & 75.0\% & 34.1\% & 74.2\% \\
    \midrule
    \multicolumn{5}{l}{\textit{Dream-7B}} \\
    0.60 & 98.8\% & 63.2\% & 85.3\% & 56.9\% \\
    0.70 & 97.9\% & 69.1\% & 66.3\% & 62.5\% \\
    0.80 & 96.7\% & 72.4\% & 36.0\% & 73.5\% \\
    0.90 & 87.3\% & 74.6\% & 6.2\% & 74.8\% \\
    \bottomrule
  \end{tabular}
  \caption{Selective generation results across thresholds. ``Saved'' is the fraction of denoising steps avoided by skipping instances predicted to fail. ``Clf.\ acc'' is the fraction of instances for which the probe correctly predicts functional vs.\ non-functional (same metric as Table~\ref{tab:early_exit}).}
  \label{tab:early_exit_full}
\end{table}

\section{Length Control Probe}
\label{sec:appendix_length}

We train a control probe (same pipeline) to predict binary output length (above/below median) instead of correctness.
To control for the confound, we re-train the correctness probe on a length-matched subsample where functional and non-functional instances have identical length distributions (binned matching, 10 quantile bins).

\vspace{0.5em}
\begin{center}
  \small
  \captionof{table}{Best AUC for correctness, length control, and length-matched correctness probes. After controlling for length, correctness AUC drops modestly but remains well above chance.}
  \label{tab:length_control}
  \vspace{0.3em}
  \begin{tabular}{lccc}
    \toprule
    & \textbf{Corr.} & \textbf{Length} & \textbf{Corr.\ (matched)} \\
    \midrule
    \multicolumn{4}{l}{\textit{JSON Schema}} \\
    ~~LLaDA & 0.85 & 0.98 & 0.81 \\
    ~~Dream & 0.84 & 0.99 & 0.75 \\
    \midrule
    \multicolumn{4}{l}{\textit{GSM8K}} \\
    ~~LLaDA & 0.79 & 0.86 & 0.75 \\
    ~~Dream & 0.82 & 0.87 & 0.80 \\
    \bottomrule
  \end{tabular}
\end{center}

\section{Spatial Pooling Ablation}
\label{sec:appendix_region}

We pool the 4 extracted regions into 1 (global average), 2 (pairwise average), or 4 (concatenation) groups and re-run the probe.

\vspace{0.5em}
\begin{center}
  \small
  \captionof{table}{Best AUC by number of position regions. Spread is less than 0.02 within each configuration.}
  \label{tab:region_ablation}
  \vspace{0.3em}
  \begin{tabular}{lccc}
    \toprule
    & \textbf{1 region} & \textbf{2 regions} & \textbf{4 regions} \\
    \midrule
    \multicolumn{4}{l}{\textit{JSON Schema}} \\
    ~~LLaDA & 0.848 & 0.854 & 0.840 \\
    ~~Dream & 0.846 & 0.828 & 0.847 \\
    \midrule
    \multicolumn{4}{l}{\textit{GSM8K}} \\
    ~~LLaDA & 0.784 & 0.788 & 0.784 \\
    ~~Dream & 0.818 & 0.814 & 0.812 \\
    \bottomrule
  \end{tabular}
\end{center}

\section{AR Baseline Probe}
\label{sec:appendix_ar}

We probe Qwen-2.5-7B-Instruct, the AR base model of Dream, using the last prompt token hidden state (one forward pass, no generation).
Sub-A uses dLLM functional correctness labels; Sub-B generates with Qwen (greedy) and uses its own correctness labels.

\vspace{0.5em}
\begin{center}
  \small
  \captionof{table}{AR vs.\ dLLM step-0 probe AUC. The AR probe uses only the last prompt token; dLLM probes use mean-pooled generation region hidden states at step~0. Sub-A predicts dLLM labels; Sub-B predicts AR's own labels.}
  \label{tab:ar_probe}
  \vspace{0.3em}
  \begin{tabular}{lcc}
    \toprule
    & \textbf{JSON Schema} & \textbf{GSM8K} \\
    \midrule
    Qwen Sub-A (dLLM labels) & 0.74 & 0.69 \\
    Qwen Sub-B (AR labels) & 0.77 & 0.77 \\
    \midrule
    LLaDA step-0 & 0.78 & 0.71 \\
    Dream step-0 & 0.80 & 0.71 \\
    \bottomrule
  \end{tabular}
\end{center}

\end{document}
